\ysrc{91}
出于若干理由（例如见猜想 4.24），把范畴 $\Rep(G(F),\Lambda)$ 稍作修改、迫使 $\delta_K$ 对所有开紧子群都成为紧对象，是方便的\tn{底本此处作「所有开紧子群 $\Lambda$」，按文意应为 $K$。}。具体地说，令
\[
\Rep_{\mathrm{f.g.}}(G(F),\Lambda)\subset\Rep(G(F),\Lambda)
\]
为这些 $\delta_K$ 在有限余极限与收缩（retract）下生成的全子范畴，并令
\[
\Ind\Rep_{\mathrm{f.g.}}(G(F),\Lambda)
\]
为其 ind-完备化。由于每个 $\delta_K$ 都是 $\Lambda$-平坦的，更换系数环时有自然的等价
\[
\Rep_{\mathrm{f.g.}}(G(F),\Lambda)\otimes_\Lambda\Lambda'=\Rep_{\mathrm{f.g.}}(G(F),\Lambda').
\]
不言而喻地，对任意开紧子群 $K\subset G(F)$，$\delta_K$ 在 $\Ind\Rep_{\mathrm{f.g.}}(G(F),\Lambda)$ 中是紧的，并且存在保持余极限的函子
\[
\Ind\Rep_{\mathrm{f.g.}}(G(F),\Lambda)\to\Rep(G,\Lambda).
\]
若 $\Lambda$ 是特征零的域，这是等价，因为由 Bernstein 的一个结果，$\Rep(G,\Lambda)^{\heartsuit}$ 具有有限的整体上同调维数。一般地，当 $\Lambda$ 是正则诺特环时，$\Ind\Rep_{\mathrm{f.g.}}(G(F),\Lambda)$ 上有自然的 $t$-结构，而上述函子限制到（关于该自然 $t$-结构的）下有界子范畴上给出等价 $\Ind\Rep_{\mathrm{f.g.}}(G(F),\Lambda)^+\cong\Rep(G,\Lambda)^+$。更详细的讨论见 [88, §3.3.3]。

对开紧子群 $K\subset G(F)$，定义相应的以 $\Lambda$ 为系数的导出 Hecke 代数
\[
H_{G,K,\Lambda}:=(\End(\delta_K))^{\mathrm{op}}.
\]
于是 $H_{G,K,\Lambda}$ 是 $\mathrm{Alg}(\Mod_\Lambda)$ 中的对象，即一个 $E_1$-代数。若 $G$ 或 $\Lambda$ 在上下文中清楚，有时省略下标中的 $G$ 或 $\Lambda$。注意其零次上同调
\[
H^0H_K\cong C_c(K\backslash G(F)/K,\Lambda)
\]
就是通常的以 $\Lambda$ 为系数的 Hecke 代数，其代数结构由卷积积给出。此外，作为 $\Lambda$-模，
\[
H_K\cong\bigoplus_{g\in K\backslash G/K}C^*(K\cap gKg^{-1},\Lambda),
\]
其中右端是 $K\cap gKg^{-1}$ 以平凡系数 $\Lambda$ 的（投射有限）群上同调。特别地，若 $K$ 是 $\Lambda$-容许的，则 $H_{G,K,\Lambda}$ 集中在零次上同调。

\begin{rem}{4.3}
通过选取 $G(F)$ 上的不变 Haar 测度、把一个（从而每个）pro-$p$ Iwahori 子群的体积取为 $1$，可以定义 $G(F)$ 的通常的完全 Hecke 代数 $H_G$。即其底空间为
\[
\delta_{\{1\}}\simeq C_c^\infty(G(F),\Lambda),
\]
乘法由通常的卷积给出。若 $K$ 是 $\Lambda$-容许的，其体积 $\mathrm{vol}(K)$ 在 $\Lambda$ 中可逆，因此照常有 $H_G$ 的幂等元 $e_K=\frac{1}{\mathrm{vol}(K)}\mathrm{ch}_\Lambda$，其中 $\mathrm{ch}_\Lambda$ 是 $\Lambda$ 的特征函数\tn{底本作「$\Lambda$ 的特征函数」，按文意应为 $K$ 的特征函数。}。范畴 $\Rep(G(F),\Lambda)^{\heartsuit}$ 与非退化 $H_G$-模范畴之间有范畴等价。作为左 $H_G$-模有 $\delta_K\simeq H_Ge_K$，且 $H_{G,K}\simeq e_KH_Ge_K$。
\end{rem}

\ysrc{92}
设 $\Mod_{H_K}$ 表示左 $H_K$-模的 $\infty$-范畴。由一般的范畴废话可知，存在一对伴随函子
\[
\delta_K\otimes_{H_K}(-):\Mod_{H_K}\rightleftarrows\Rep(G(F),\Lambda):\Hom(\delta_K,-).
\]
若 $K$ 是 $\Lambda$-容许的，则 $W\mapsto\delta_K\otimes_{H_K}W$ 是全忠实的。（若把 $\Rep(G(F),\Lambda)$ 换成 $\Ind\Rep_{\mathrm{f.g.}}(G(F),\Lambda)$，则对任意 $K$ 都全忠实。）

对 $G(F)$ 的两个开紧子群 $K_1$ 与 $K_2$，有 $(H_{K_2}\times H_{K_1})$-双模
\[
K_1H_{K_2}:=\Hom(\delta_{K_1},\delta_{K_2}).
\]
其零次上同调为
\[
H^0(K_1H_{K_2})\cong C_c(G(F)/K_2)^{K_1}=:C_c(K_1\backslash G(F)/K_2),
\]
即 $G(F)$ 上 $(K_1\times K_2)$-不变的紧支撑函数空间。若 $K_1$ 与 $K_2$ 之中有一个是 $\Lambda$-容许的，则 $K_1H_{K_2}=H^0(K_1H_{K_2})$。

不言而喻地，在上述等同下，把 $\mathrm{ch}_{K_1}\in\delta_{K_1}$ 送到 $\mathrm{ch}_{K_1K_2}\in\delta_{K_2}$ 的映射 $\iota_{K_1,K_2}:\delta_{K_1}\to\delta_{K_2}$ 对应于 $\mathrm{ch}_{K_1K_2}\in C_c(K_1\backslash G(F)/K_2)$。另一方面，
\[
\mathrm{Av}_{K_1,K_2}:\delta_{K_1}\to\delta_{K_2},\qquad
(\mathrm{Av}_{K_1,K_2}f)(g)=\int_{K_2}f(gk)\,dk
\]
对应于 $\mathrm{vol}(K_2)\mathrm{ch}_{K_1K_2}$。

不言而喻地，存在 $G(F)$-模同态
\begin{equation}
\tag{4.1}
\delta_{K_1}\otimes_{H_{K_1}}K_1H_{K_2}\to\delta_{K_2}.
\end{equation}
若 $K_1\subset K_2$ 且 $K_2$ 是 $\Lambda$-容许的开紧子群（于是 $K_1$ 也是），则 (4.1) 是同构。但一般情形未必如此。

\begin{exam}{4.4}
令 $G=\mathrm{SL}_2$，$K_2=K=\mathrm{SL}_2(\mathcal{O}_F)$，$K_1=I$ 为标准 Iwahori 子群。令 $\Lambda=\mathbb{F}_\ell$，其中 $\ell>2$ 且 $\ell\mid p+1$。则 $I$ 是 $\Lambda$-容许的，而 $K$ 不是。此时 (4.1) 不是同构。事实上，$\delta_I\otimes_{H_I}{}_IH_K$ 甚至不集中在零次。
\end{exam}

下面简要回顾 Whittaker 模。设 $G$ 在 $F$ 上拟分裂，且 $\Lambda$ 是诺特的 $\Z[1/p]$-代数、包含所有 $p$ 次单位根（例如 $\Lambda=W(\overline{\mathbb{F}}_\ell)$）。$G$ 的一个 Whittaker 数据由一个 $F$-有理 Borel 子群的幂幺根基 $U$ 与一个非退化特征 $\psi:U(F)\to(U/[U,U])(F)\to\Lambda^\times$ 组成。给定 Whittaker 数据 $(U,\psi)$，令
\[
\mathrm{Whit}_{U,\psi}:=c\text{-}\mathrm{ind}_{U(F)}^{G(F)}\psi\in\Rep(G(F),\Lambda)^{\heartsuit}
\]
为相应的 Whittaker 模。注意 $\mathrm{Whit}_{U,\psi}$ 作为 $G(F)$-模不是有限生成的。然而它可以写成 $\Rep(G(F),\Lambda)^{\heartsuit}$ 中有限生成投射对象的滤过余极限（[63, Prop. 3]）。

在本小节末尾，回顾 $\Rep(G(F),\Lambda)$ 的一些内部对称性。首先回忆，每个拓扑群自同构 $c:G(F)\to G(F)$ 都诱导范畴的自等价 $c:\Rep(G(F),\Lambda)^{\heartsuit}\to\Rep(G(F),\Lambda)^{\heartsuit}$。即，若 $V$ 是 $G(F)$ 的光滑表示，定义新表示 ${}^cV$，使 ${}^cV=V$ 作为 $\Lambda$-模，但带有新的 $G(F)$-作用 $G(F)\times{}^cV\to{}^cV$，$(g,v)\mapsto c^{-1}(g)v$。若 $K$ 是开紧子群，则有典范同构
\[
{}^c\delta_K\simeq\delta_{c(K)},\qquad f\mapsto{}^cf,\quad ({}^cf)(x)=f(c^{-1}(x)).
\]

\ysrc{93}
把这一形式体系应用于 $G_{\mathrm{ad}}(F)$ 通过内自同构在 $G(F)$ 上的作用，得到 $G_{\mathrm{ad}}(F)$ 在 $\Rep(G(F),\Lambda)^{\heartsuit}$ 上的作用。注意若 $c=c_h$ 由元素 $h\in G(F)$ 的共轭给出，则有典范同构 ${}^{c_h}V\simeq V$，$v\mapsto hv$。由此 $G_{\mathrm{ad}}(F)$ 在 $\Rep(G(F),\Lambda)^{\heartsuit}$ 上的作用穿过 Picard 群胚
\[
\mathbf{Tor}^0_{Z_G}:=G_{\mathrm{ad}}(F)/G(F)
\]
的作用，后者延拓为作用
\begin{equation}
\tag{4.2}
\mathbf{Tor}^0_{Z_G}\times\Rep(G(F),\Lambda)\to\Rep(G(F),\Lambda).
\end{equation}
注意 $\mathbf{Tor}^0_{Z_G}$ 可以等同于 $F$ 上 $Z_G$-挠子中那些使诱导的 $G$-挠子平凡的挠子构成的 Picard 群胚。特别地，$\mathbf{Tor}^0_{Z_G}$ 的同构类群是
\[
E_G:=\pi_0\mathbf{Tor}^0_{Z_G}=G_{\mathrm{ad}}(F)/(G(F)/Z_G(F))\cong\ker\big(H^1(F,Z_G)\to H^1(F,G)\big),
\]
而 $\mathbf{Tor}^0_{Z_G}$ 中任一对象的自同构群是 $Z_G(F)$。

还有所谓的 $\Rep_{\mathrm{f.g.}}(G(F),\Lambda)$ 的上同调对偶函子 $\mathbb{D}^{\mathrm{coh}}$：
\begin{equation}
\tag{4.3}
\mathbb{D}^{\mathrm{coh}}:\Rep_{\mathrm{f.g.}}(G(F),\Lambda)\to\Rep_{\mathrm{f.g.}}(G(F),\Lambda)^{\mathrm{op}},\qquad
V\mapsto\Hom_{G(F)}(V,H_G),
\end{equation}
其中 $H_G=C_c^\infty(G(F),\Lambda)$ 是完全 Hecke 代数，看作其自身上的双模。

\subsection{群胚 $\mathbf{W}_G$ 与 $\mathbf{TS}_G$}
在一般约化群的局部 Langlands 对应的标准表述中，必须作若干辅助选择。我们的表述同样如此，这一点将在后文讨论。本小节说明如何仔细地选取这些辅助数据。与现有文献相比，我们将引入一些群胚来记录这些数据的自同构。主要关心满足条件 $H^1(F,Z_G)=0$ 的拟分裂群（例如 $G=\mathrm{GL}_n$）的读者可以基本跳过本小节。

设 $\mathrm{Pin}_G$ 表示 $G$ 的固定数据（pinning）簇。即，对经典的 $F$-代数 $A$，$\mathrm{Pin}_G(A)$ 由三元组 $(B_A,T_A,e_A)$ 组成，其中 $B_A\subset G_A$ 是 Borel 子群，$T_A\subset B_A$ 是极大环面，$e_A:U_A\to\mathbb{G}_a$ 是 $U_A$（$B_A$ 的幂幺根基）上的同态，使得经过某个使 $G_{A'}$ 分裂的 étale 覆盖 $A\to A'$ 之后，$e_A$ 限制到每个对应于单根（相对于 $(B_A,T_A)$）的根子群上给出同构 $U_\alpha\to\mathbb{G}_a$。注意 $\mathrm{Pin}_G$ 实际上是 $G_{\mathrm{ad}}$-挠子。其上同调类 $\alpha\in H^1(F,G_{\mathrm{ad}})$ 对应于 $G$ 的拟分裂内型。特别地，$\mathrm{Pin}_G$ 有有理点当且仅当 $G$ 拟分裂；此时 $\mathrm{Pin}_G(F)$ 是 $G_{\mathrm{ad}}(F)$-挠子。于是若 $G$ 拟分裂，可以定义商群胚
\begin{equation}
\tag{4.4}
\mathbf{W}_G:=\mathrm{Pin}_G(F)/G(F).
\end{equation}
注意它是 $\mathbf{Tor}^0_{Z_G}$-挠子，故其同构类集合 $W_G=\pi_0\mathbf{W}_G$ 是 $E_G$-挠子。

本小节的第一个目标是以 $G_{\mathrm{ad}}(F)$-等变的方式，把若干对象典范地联系到 $(B,T,e)\in\mathrm{Pin}_G(F)$。

\ysrc{94}
首先，若选取非平凡加性特征 $\psi_0:F\to\Lambda^\times$（因而特别地我们假设 $\Lambda$ 包含足够多的 $p$ 次单位根），则有定义良好的 $G_{\mathrm{ad}}(F)$-等变映射，从 $\mathrm{Pin}_G(F)$ 到 $G$ 的 Whittaker 数据集合，把 $(B,T,e)$ 送到
$\big(U,\ \psi:U(F)\xrightarrow{e}F\xrightarrow{\psi_0}\Lambda^\times\big)$；它在 $W_G$ 与 Whittaker 数据的 $G(F)$-共轭类集合之间诱导双射。于是有定义良好的 $\mathbf{Tor}^0_{Z_G}$-等变函子
\begin{equation}
\tag{4.5}
\mathfrak{W}_{\psi_0}:\mathbf{W}_G\to\Rep(G(F),\Lambda),\qquad (B,T,e)\mapsto\mathrm{Whit}_{U,\psi}.
\end{equation}

\begin{rem}{4.5}
由于 $\mathfrak{W}_{\psi_0}$ 在猜想表述中要用到，这里简要讨论它如何依赖于 $\psi_0$ 的选取。给定 $\psi_0$ 与 $\psi_0'$，存在唯一的 $a\in F^\times$ 使 $\psi_0'(-)=\psi_0(a-)$。给定固定数据 $(B,T,e)$，若 $a$ 在映射 $F^\times\xrightarrow{\hat\rho}T_{\mathrm{ad}}(F)\to H^1(F,Z_G)$ 下的像平凡，则两个 Whittaker 模 $\mathrm{Whit}_{U,\psi_0e}$ 与 $\mathrm{Whit}_{U,\psi_0'e}$ 同构，其中 $\hat\rho$ 是 $G$ 的正余根之和的一半。因此若 $H^1(F,Z_G)$ 平凡，则 $\mathfrak{W}_{\psi_0}$ 与 $\psi_0$ 的选取无关（至同构）。一般地，它至多依赖 $a$ 在 $F^\times/(F^\times)^2$ 中的像。此外，在局部情形总可假设 $\psi_0$ 的导子是 $\mathcal{O}_F$（即 $\psi_0|_{\mathcal{O}_F}=1$ 但 $\psi_0|_{\varpi_F^{-1}\mathcal{O}_F}\neq1$），从而把歧义缩小到 $\kappa_F^\times/(\kappa_F^\times)^2$。我们还想指出，应当可以在不涉及 $\psi_0$ 的选取（并允许 $\Lambda$ 不包含足够多 $p$ 次单位根）的情况下更典范地表述一切，不过我们选择不这样做。
\end{rem}

下一步，我们构造从 $\mathrm{Pin}_G(F)$ 到由 Iwahori 子群 $I$ 与 $G(F)$ 的特殊 parahoric $K$ 组成的对 $(I\subset K)$ 的 $G_{\mathrm{ad}}(F)$-等变映射。$\breve F$ 仍表示 $F$ 的极大非分歧扩域 $F^{\mathrm{ur}}$ 的完备化，并令
\begin{equation}
\tag{4.6}
\kappa_G:G(\breve F)\to\mathbb{X}^\bullet(Z_{\hat G}^{I_F})
\end{equation}
为 Kottwitz 映射（[47, §7]）。选取 $\Z$ 上带固定数据的 Chevalley 群 $(H,B_H,T_H,e_H)$ 以及同构
$\eta:(H,B_H,T_H,e_H)_{\breve F}\simeq(G,B,T,e)_{\breve F}$。则
\[
K=\eta\big(H(\mathcal{O}_{\breve F})\big)^{\Gamma_{\breve F/F}}\cap\ker\kappa_G
\]
（交在 $G(F)$ 中取）是特殊 parahoric，且与 $(H,B_H,T_H,e_H,\eta)$ 的选取无关。设 $S\subset T$ 是极大 $F$-分裂环面。利用对应于 $\Lambda$ 的特殊顶点 $x\in A(G,S)$\tn{底本此处作「对应于 $\Lambda$ 的顶点」，按文意应为 $K$；下页一处同此。}，可以把公寓 $A(G,S)$（在 $G$ 的 Bruhat–Tits 建筑中）等同于由 $S$ 的余权格张成的实向量空间。则 $I$ 是唯一的 Iwahori 子群，其对应的 alcove $\mathbf{a}$ 包含 $x$，且含于 $B$ 所确定的有限 Weyl 房内。

\begin{rem}{4.6}
上述构造的特殊 parahoric $K$ 是绝对特殊的，即其在 $G$ 的 Bruhat–Tits 建筑中对应的顶点 $x$ 对每个有限可分扩张 $F'/F$ 仍是特殊的（另见 [13, §3] 末尾）。

在 [84, §6] 中引入了密切相关的记号：$G$ 的特殊 parahoric 若对应顶点对每个非分歧扩张 $F'/F$ 仍是特殊的，则称为很特殊的（very special）。显然绝对特殊的 parahoric 是很特殊的，因此仅当 $G$ 拟分裂时才存在（由 loc. cit. 的引理 6.1）。另一方面，对拟分裂的 $G$，上述构造给出绝对特殊 parahoric 的一个 $G_{\mathrm{ad}}(F)$-共轭类。事实上，借助以下事实，该构造给出全部绝对特殊 parahoric。
\end{rem}

\begin{lem}{4.7}
所有绝对特殊 parahoric 在 $G_{\mathrm{ad}}(F)$-作用下共轭。
\end{lem}

该引理推广了熟知的事实：非分歧群 $G$ 中所有超特殊 parahoric 在 $G_{\mathrm{ad}}(F)$ 下共轭。为证明该引理，可假设 $G=G_{\mathrm{ad}}$ 且拟分裂、绝对单。

\ysrc{95}
于是由分类容易推出该引理。然而注意，该引理对很特殊的 parahoric 不成立。事实上，对奇分歧酉群 $\mathrm{U}_{2m+1}$（例如设 $\operatorname{char}\kappa_F\neq2$），有两类很特殊的 parahoric，其一约化商为 $\mathrm{SO}_{2m+1}$，另一约化商为 $\mathrm{Sp}_{2m}$（例如见 [84]）。只有前者是绝对特殊的。

设 $\widetilde W=N_G(T)(\breve F)/\ker\kappa_T$ 为 $G_{\breve F}$ 关于 $T_{\breve F}$ 的 Iwahori–Weyl 群，它适合短正合列 $1\to\mathbb{X}^\bullet(\hat T^{I_F})\to\widetilde W\to W_0\to1$，其中 $W_0$ 如前是 $G_{\breve F}$ 的有限 Weyl 群。由于对应于 $\Lambda$ 的顶点 $x$ 对 $G_{\breve F}$ 仍是特殊的，它给出上述序列的一个分裂，于是可以写
\begin{equation}
\tag{4.7}
\widetilde W=\mathbb{X}^\bullet(\hat T^{I_F})\rtimes W_0.
\end{equation}
alcove $\mathbf{a}$ 对 $G_{\breve F}$ 仍是 alcove（对应一个 Iwahori 子群 $\breve I\subset G(\breve F)$），并确定子群
\begin{equation}
\tag{4.8}
\Omega\cong N_{G(\breve F)}(\breve I)/\breve I\subset\widetilde W
\end{equation}
它固定该 alcove。众所周知，Kottwitz 映射 (4.6) 诱导同构 $\Omega\simeq\mathbb{X}^\bullet(Z_{\hat G}^{I_F})$。因此每个 $\gamma\in\mathbb{X}^\bullet(Z_{\hat G}^{I_F})$ 可唯一地写成
\begin{equation}
\tag{4.9}
\gamma=\lambda_\gamma w_\gamma,\qquad \lambda_\gamma\in\mathbb{X}^\bullet(\hat T^{I_F}),\quad w_\gamma\in W_0.
\end{equation}

设 $H_I$ 是 $I$ 的 Iwahori Hecke 代数。注意 $I\cap T(F)$ 是 $T$ 的 Iwahori 子群，故有相应的 Iwahori Hecke 代数 $H_{T,I}$。类似地有 pro-$p$ Iwahori Hecke 代数 $H_{I(1)}$ 与 $H_{T,I(1)}$。已知\footnote{感谢 Vigneras 指出这一点。}作为 $G(F)$-表示，
\[
\delta_I\cong\Ind_{B(F)}^{G(F)}\delta_{T,I},\qquad
\delta_{I(1)}\cong\Ind_{B(F)}^{G(F)}\delta_{T,I(1)},
\]
其中 $\delta_{T,I}$ 与 $\delta_{T,I(1)}$ 是 $T(F)$ 由其 Iwahori 子群与 pro-$p$ Iwahori 子群紧诱导的表示。（这些同构在 $\Lambda=\C$ 时大概是熟知的；对一般使 $p$ 可逆的 $\Lambda$，它们隐含在 [18, 3.6, 6.2, 6.3] 中。）由此有代数的典范映射
\begin{equation}
\tag{4.10}
H_{T,I}\to H_I,\qquad H_{T,I(1)}\to H_{I(1)},
\end{equation}
它们在取 $H^0$ 之后是单射。它们正是 Bernstein 构造的（pro-$p$）Iwahori Hecke 代数的交换子代数。另一方面，由
\[
\delta_I=c\text{-}\mathrm{ind}_K^{G(F)}c\text{-}\mathrm{ind}_I^K\mathbf{1},\qquad
\delta_{I(1)}=c\text{-}\mathrm{ind}_K^{G(F)}c\text{-}\mathrm{ind}_{I(1)}^K\mathbf{1}
\]
得到典范映射
\begin{equation}
\tag{4.11}
H_f:=\End_K(c\text{-}\mathrm{ind}_I^K\mathbf{1})^{\mathrm{op}}\to H_I,\qquad
H_{f,(1)}:=\End_K(c\text{-}\mathrm{ind}_{I(1)}^K\mathbf{1})^{\mathrm{op}}\to H_{I(1)}.
\end{equation}

\begin{rem}{4.8}
我们指出，Iwahori–Weyl 群与分解 (4.7)、(4.9)，以及（pro-$p$）Iwahori Hecke 代数与 (4.10)、(4.11)，都典范地附着于 $W_G$ 中的一个元素。事实上，若 $(B_1,T_1,e_1)$ 到 $(B_2,T_2,e_1)$\tn{底本第二个三元组中的 $e_1$ 疑为 $e_2$ 之误。} 是同一 $G(F)$-共轭类中的两组固定数据，则把一个共轭到另一个的 $g\in G(F)$ 的选取诱导这些数据之间的同构，且同构实际上与 $g$ 的选取无关。
\end{rem}

\ysrc{96}
\begin{rem}{4.9}
一个有趣的问题是：由 (4.10) 与 (4.11) 诱导的 $\Lambda$-模映射 $H_{T,I}\otimes_\Lambda H_f\to H_I$ 是否为同构。已知取 $H^0$ 之后确实如此。
\end{rem}

下面再提一个结果。

\begin{prop}{4.10}
固定导子为 $\mathcal{O}_F$ 的加性特征 $\psi_0:F\to\Lambda^\times$。对应 $(B,T,e)\mapsto(U,\psi)$ 与 $(B,T,e)\mapsto(I\subset K)$ 诱导定义良好的 $E_G$-等变映射 $(U,\psi)\mapsto(I\subset K)$，从 Whittaker 数据的 $G(F)$-共轭类集合到由绝对特殊 parahoric $K$ 与 Iwahori 子群 $I\subset K$ 组成的对的 $G(F)$-共轭类集合。该对应与 $\psi_0$ 的选取无关。若 $(U,\psi)$ 映到 $(I\subset K)$，则 $\mathrm{Whit}_{U,\psi}^K$ 是秩为 $1$ 的自由 $H^0H_K$-模（即 Casselman–Shalika 公式 [13]），而 $\mathrm{Whit}_{U,\psi}^I\simeq M_{\mathrm{asp}}$ 是 $H^0H_I$ 的反球模（antispherical module）（即由 $H_f\subset H^0H_I$ 的符号表示诱导的表示）。
\end{prop}

关于拟分裂群的讨论到此结束。为讨论未必分裂的约化群，先注意：从几何观点看，考虑分类 $\mathrm{Pin}_G$ 到 $G$-挠子的提升的群胚 $(\mathrm{Pin}_G/G)(F)$ 更为自然，它把 $\mathrm{Pin}_G(F)/G(F)$ 作为子群胚包含在内。注意 $(\mathrm{Pin}_G/G)(F)$ 是中性带叠（neutral gerbe），更确切地说，是 $F$ 上 $Z_G$-挠子的 Picard 群胚 $\mathbf{Tor}_{Z_G}$ 下的挠子（特别地受 $\mathbf{Tor}^0_{Z_G}\subset\mathbf{Tor}_{Z_G}$ 作用）。即使 $G$ 不拟分裂、$\mathrm{Pin}_G(F)=\varnothing$，仍可考虑群胚 $(\mathrm{Pin}_G/G)(F)$，它可能非空。更确切地说，它非空当且仅当类 $\alpha\in H^1(F,G_{\mathrm{ad}})$ 可以提升为 $H^1(F,G)$ 中的类；此时它仍是 $\mathbf{Tor}_{Z_G}$-挠子。然而对许多应用而言，这个群胚仍不够大，因为 $\alpha$ 常常不能提升为 $H^1(F,G)$ 中的类。因此我们引入更大的群胚 $\mathbf{TS}_G$，它对大多数应用足够。

首先，类似于 §3.2 中引入的群胚 $\mathbf{Tor}_{Z,\mathrm{iso}F}$，令 $\mathbf{Tor}_{G,\mathrm{iso}F}$ 为由 $\breve F$ 上的 $G$-挠子 $\mathcal{E}$ 与 $G$-挠子的同构 $\varphi:\mathcal{E}\simeq\sigma^*\mathcal{E}$ 组成的对 $(\mathcal{E},\varphi)$ 的群胚。其同构类集合正是 Kottwitz 集 $B(G)$（[46,47]）。给定 $\mathbf{Tor}_{G,\mathrm{iso}F}$ 中的 $b=(\mathcal{E},\varphi)$，可定义 $F$-代数群 $G_b$，其 $A$-点（对经典 $F$-代数 $A$）构成 $(\mathcal{E}\otimes_FA,\varphi\otimes1)$ 在 $\breve F\otimes_FA$ 上的自同构群。Kottwitz 证明，在 $\breve F$ 上 $G_b$ 自然同构于 $G$ 的一个 Levi 子群。若它在 $\breve F$ 上同构于 $G$（此时 $G_b$ 自然是 $G$ 的内型），则称 $b$ 是基本的（basic）。基本 $b$ 的同构类集合记作 $B(G)_{\mathrm{bsc}}$。把 $\mathcal{E}\mapsto(\mathcal{E}\otimes_F\breve F,\varphi=1\otimes\sigma)$ 给出从 $F$ 上 $G$-挠子的范畴 $\mathbf{Tor}_G$ 到 $\mathbf{Tor}_{G,\mathrm{iso}F}$ 的全忠实嵌入，它诱导嵌入 $H^1(F,G)\subset B(G)_{\mathrm{bsc}}$。回忆 Kottwitz 的以下上同调结果。
\begin{itemize}[leftmargin=1.6em,itemsep=0.2em]
\item 对每个 $G$，映射 (4.6) 诱导映射 $\kappa_G:B(G)\to\mathbb{X}^\bullet(Z_{\hat G}^{\Gamma_F})$（仍称为 Kottwitz 映射），它限制为双射 $B(G)_{\mathrm{bsc}}\simeq\mathbb{X}^\bullet(Z_{\hat G}^{\Gamma_F})$。
\item 自然映射 $H^1(F,G)\to B(G)_{\mathrm{bsc}}$ 在 $G=G_{\mathrm{ad}}$ 为伴随型时是双射。
\end{itemize}

现在可通过嵌入 $\mathbf{Tor}_{G_{\mathrm{ad}}}\subset\mathbf{Tor}_{G_{\mathrm{ad}},\mathrm{iso}F}$ 把 $\mathrm{Pin}_G$ 看作 $\mathbf{Tor}_{G_{\mathrm{ad}},\mathrm{iso}F}$ 中的对象，并考虑 $\mathrm{Pin}_G$ 到 $\mathbf{Tor}_{G,\mathrm{iso}F}$ 中对象的提升所构成的群胚 $\mathbf{TS}_G$。具体地说，$\mathbf{TS}_G$ 的对象由 $t=(b,B,T,e)$ 组成，其中 $b=(\mathcal{E},\varphi)\in\mathbf{Tor}_{G,\mathrm{iso}F}$ 是基本的，$(B,T,e)$ 是 $G_b$ 的固定数据。

\ysrc{97}
对象 $t$ 与 $t'$ 之间的态射是 $\mathbf{Tor}_{G,\mathrm{iso}F}$ 中 $b$ 与 $b'$ 之间的同构，它诱导同构 $(G_b,B,T,e)\simeq(G_{b'},B',T',e')$。该群胚非空当且仅当 $\alpha$ 可提升为 $\mathbb{X}^\bullet(Z_{\hat G}^{\Gamma_F})$ 中的元素。这仍未必总能做到。例如若 $G=D^{\mathrm{Nm}=1}$ 是 $F$ 上四元数代数 $D$ 中约化范数为 $1$ 的元素群，则这样的延拓不存在。然而，若 $G$ 拟分裂或 $G$ 的中心连通，则这样的提升总存在；此时 $\mathbb{X}^\bullet(Z_{\hat G}^{\Gamma_F})\to\mathbb{X}^\bullet(Z_{\hat G_{sc}}^{\Gamma_F})$ 是满射（回忆 $\hat G_{sc}$ 表示 $G_{\mathrm{ad}}$ 的对偶群，即 $\hat G$ 的导群的单连通覆盖）。若 $\mathbf{TS}_G$ 非空，则它是 $\mathbf{Tor}_{Z_G,\mathrm{iso}F}$ 下的挠子（故其同构类集合 $\pi_0\mathbf{TS}_G$ 是 $B(Z_G)$ 下的挠子）。注意若 $G$ 拟分裂，则 $\mathbf{W}_G\subset\mathbf{TS}_G$ 且
\[
\mathbf{TS}_G=\mathbf{W}_G\times_{\mathbf{Tor}^0_{Z_G}}\mathbf{Tor}_{Z_G,\mathrm{iso}F}.
\]

现在固定 $t\in\mathbf{TS}_G$，并把 $(G_b,B,T,e)$ 记作 $(G^*,B^*,T^*,e^*)$。可以典范地把对偶群 $\hat G$ 与对偶群 $\hat G^*$ 等同。附着于 $(G^*,B^*,T^*,e^*)$ 的有各种对象，如 Iwahori–Weyl 群 $\widetilde W^*=\mathbb{X}^\bullet(\hat T^{I_F})\rtimes W_0^*$ 与 Iwahori–Hecke 代数 $H_{I^*}$。$b$ 的类是元素 $\beta\in B(G)_{\mathrm{bsc}}\simeq\mathbb{X}^\bullet(Z_{\hat G}^{\Gamma_F})$，它是类 $\alpha\in\mathbb{X}^\bullet(Z_{\hat G_{sc}}^{\Gamma_F})$ 的提升。此外，对 $-\beta$ 沿 $\mathbb{X}^\bullet(Z_{\hat G}^{I_F})\twoheadrightarrow\mathbb{X}^\bullet(Z_{\hat G}^{\Gamma_F})$ 的每个提升 $\gamma$，得到 $G(F)$ 的一个典范定义的 Iwahori–Hecke 代数 $H_{I_\gamma}$。即，若进一步把 $\gamma$ 沿 $N_{G^*(\breve F)}(\breve I^*)\twoheadrightarrow\Omega^*\cong\mathbb{X}^\bullet(Z_{\hat G}^{I_F})$ 提升为元素 $\tilde\gamma$，则得到 $G(F)$ 的 Iwahori 子群 $I_{\tilde\gamma}$。事实上，利用 $\mathcal{E}$ 可以等同
\[
G(F)\cong\{g\in G^*(\breve F)\mid\tilde\gamma\sigma(g)\tilde\gamma^{-1}=g\}.
\]
则 $I_{\tilde\gamma}=\{g\in I^*_{\breve F}\mid\tilde\gamma\sigma(g)\tilde\gamma^{-1}=g\}$。相应的 Iwahori–Hecke 代数只依赖 $\gamma$，故记作 $H_{I_\gamma}$。

\subsection{导出 Satake 同构}
固定 $\iota:\Gamma_q\to\Gamma_F^t$，于是有 $\Z[1/p]$ 上的叠 $\Loc_{{}^cG,F,\iota}$。本小节假设 $G$ 非分歧。则有 $\Loc^{ur}_{{}^cG,F}\subset\Loc^{tame}_{{}^cG,F,\iota}$。设 $\Lambda$ 是正则诺特 $\Z[1/p]$-代数。我们用同样的记号表示这些叠到 $\Lambda$ 的基变换。我们的第一个猜想可以看作导出 Satake 同构。\footnote{作者在 Oaxaca 的「Modularity and Moduli Spaces」会议上想到这个猜想，其灵感来自 Emerton 希望「看见」导出 Hecke 代数在模曲线（以及一般 Shimura 簇）上同调上的作用，并受到 Feng 关于谱 Hecke 代数 [28] 的结果的鼓励。讨论见注 4.61。}

\begin{conj}{4.11}
对每个超特殊子群 $K$，存在 $\Lambda$-代数的自然同构
\[
H_K\cong\Big(\widehat{\End}_{\OO_{\Loc_{{}^cG,F,\iota}}}(\OO_{\Loc^{ur}_{{}^cG,F}})\Big)^{\mathrm{op}},
\]
它在取 $H^0$ 之后化为经典的 Satake 同构：
\[
C_c(K\backslash G(F)/K,\Lambda)\cong H^0H_K\cong H^0\End_{\OO_{\Loc_{{}^cG,F,\iota}}}(\OO_{\Loc^{ur}_{{}^cG,F}})
\cong H^0\Gamma\big(\Loc^{ur}_{{}^cG,F},\OO_{\Loc^{ur}_{{}^cG,F}}\big).
\]
此外，该同构与命题 3.12 中关于不同 $\iota$ 的同构相容。
\end{conj}

由于 $\Loc^{tame}_{{}^cG,F,\iota}$ 是 $\Loc_{{}^cG,F,\iota}$ 中的开闭子叠，可以把上述猜想中的 $\OO_{\Loc_{{}^cG,F,\iota}}$ 换成 $\OO_{\Loc^{tame}_{{}^cG,F,\iota}}$。

\ysrc{98}
\begin{rem}{4.12}
\begin{enumerate}[label=(\arabic*),leftmargin=1.8em,itemsep=0.15em]
\item 注意即使 $\Lambda=\C$，该猜想也非平凡：它等价于说
$\End_{\OO_{\Loc_{{}^cG}}}(\OO_{\Loc^{ur}_{{}^cG}})=\End_{\OO_{\Loc^{unip}_{{}^cG}}}(\OO_{\Loc^{ur}_{{}^cG}})$ 集中在零次。这可由下文的定理 4.21 推出。但我们请读者直接对 $G=\mathrm{GL}_2$ 验证，以体会其内容。
\item 几何 Langlands 提示 $H_K$ 与 $\big(\widehat{\End}_{\OO_{\Loc_{{}^cG,F,\iota}}}(\OO_{\Loc^{ur}_{{}^cG,F}})\big)^{\mathrm{op}}$ 都带有自然的交换结构（使它们成为 $E_3$-代数）\footnote{在等特征情形看到这一点的一种可能方式，是取几何 Langlands 中相应 $E_3$-范畴的迹。}，尽管我们看不出如何直接构造这样的结构。若果真如此，人们还会期望上述猜想中的同构尊重交换结构。注意 $H_K$ 上 $E_3$-结构的存在将蕴含上同调环 $\oplus_iH^iH_K$ 是分次交换的，而这目前只在关于基环 $\Lambda$ 的某些假设下已知（[72]）。
\item 用这种方式表述 mod $p$ 的导出 Satake 同构（甚至整系数导出 Satake 同构）将是很有意思的。整系数的非导出情形出现在 [87]，其表述涉及 $\hat G$ 的 Vinberg 幺半群。
\end{enumerate}
\end{rem}

当 $G=T$ 是非分歧环面时，可以手工验证该猜想。

\begin{prop}{4.13}
猜想 4.11 对非分歧环面成立。
\end{prop}

\begin{proof}
由 (3.6) 有
\[
\End_{\OO_{\Loc^{tame}_{{}^cT,F,\iota}}}(\OO_{\Loc^{ur}_{{}^cT,F}})
\simeq\End_{({}^{cl}\mathcal{R}_{\kappa_{\breve F}^\times,\hat T})^\sigma}\OO_{\{1\}}
\otimes\Gamma(\hat T/(\sigma-1)\hat T,\OO).
\]
另一方面有典范同构 $H_K\cong C^*(T(\kappa_F),\Lambda)\otimes H^0H_K$。于是所要的同构来自经典 Satake 同构
\[
\Gamma(\hat T/(\sigma-1)\hat T,\OO)\cong H^0H_K
\]
以及典范同构（下文构造）
\begin{equation}
\tag{4.12}
\Lambda\big[({}^{cl}\mathcal{R}_{\kappa_{\breve F}^\times,\hat T})^\sigma\big]\simeq\Lambda T(\kappa_F),
\end{equation}
其中回忆左端是 $({}^{cl}\mathcal{R}_{\kappa_{\breve F}^\times,\hat T})^\sigma$ 的正则函数环，右端是 $T(\kappa_F)$ 的群环。

为构造 (4.12)，先假设 $T$ 分裂，于是 $\sigma$ 在 $\hat T$ 上平凡作用且 $\breve F=F$。则
\[
\Lambda\big[({}^{cl}\mathcal{R}_{\kappa_F^\times,\hat T})^\sigma\big]=\Lambda\big[\mathbb{X}_\bullet(T)\otimes\kappa_F^\times\big],
\]
且 $\mathbb{X}_\bullet(T)\otimes\kappa_F^\times\cong T(\kappa_F)$，其中 $\mathbb{X}_\bullet(T)$ 表示 $T$ 的余特征标格（定义在 $F$ 上）。利用范映射 $\mathrm{Res}_{\kappa_{\breve F}/\kappa_F}T_{\kappa_{\breve F}}\to T_{\kappa_F}$，一般非分歧环面的构造 (4.12) 化归分裂情形。
\end{proof}

\ysrc{99}
\subsection{凝聚 Springer 层}
本小节假设 $\breve F/F$ 驯顺分歧。我们在 $\Loc^{tame}_{{}^cG,F,\iota}$ 上定义一个（凝聚层的）复形，并讨论它的一些（猜想中的）性质。由于其定义与 Springer 层相像，我们称之为凝聚 Springer 层\footnote{这个名称来自 D. Ben-Zvi。}。和以前一样，所有叠都基变换到 $\Lambda$。

回忆态射 $\pi^{tame}:\Loc^{tame}_{{}^cB,F,\iota}\to\Loc^{tame}_{{}^cG,F,\iota}$，并记 $\pi^{unip}:\Loc^{unip}_{{}^cB,F,\iota}\to\Loc^{tame}_{{}^cG,F,\iota}$。对 $?=\mathrm{tame}$ 与 $\mathrm{unip}$，令
\[
\mathrm{CohSpr}^?_{{}^cG,F,\iota}:=\pi^?_*\OO_{\Loc^?_{{}^cB,F,\iota}}\in\Coh(\Loc^{tame}_{{}^cG,F,\iota}).
\]
再强调所有函子都是导出的。先注意 $\mathrm{CohSpr}^?_{{}^cG,F,\iota}$ 的以下性质。

\begin{prop}{4.14}
（凝聚层）复形 $\mathrm{CohSpr}^?_{{}^cG,F,\iota}$ 关于 $\Loc^{tame}_{{}^cG,F,\iota}$ 上的 Grothendieck–Serre 对偶是自对偶的。
\end{prop}

\begin{proof}
由命题 2.33 与注 2.34，$\Loc^{tame}_{{}^cB,F,\iota}$ 是拟光滑的且对偶复形平凡。对 $\Loc^{unip}_{{}^cB,F,\iota}$ 同样如此。因此可以在 $\mathrm{CohSpr}^?$ 的定义中用 $\Loc^?_{{}^cB,F,\iota}$ 的对偶复形 $\omega_{\Loc^?_{{}^cB,F,\iota}}$ 替换 $\OO_{\Loc^?_{{}^cB,F,\iota}}$。断言随即成立，因为 Grothendieck–Serre 对偶与真前推交换。
\end{proof}

我们在 §4.7 中的猜想提示，凝聚 Springer 层与自守提升定理中的拼接模有关。正如 Emerton 向我们解释的，拼接模应当是（经典）形变环上的（通常）极大 Cohen–Macaulay 模。这引导我们作出以下猜想（当 $G$ 分裂且 $\Lambda=\C$ 时另见 [5, 3.15]）。

\begin{conj}{4.15}
复形 $\mathrm{CohSpr}^?_{{}^cG,F,\iota}$ 属于凝聚层阿贝尔范畴 $\Coh(\Loc^{tame}_{{}^cG,F,\iota})^{\heartsuit}$。
\end{conj}

\begin{cor}{4.16}
假设猜想 4.15，则 $\mathrm{CohSpr}^?_{{}^cG,F,\iota}$ 是 $\Loc^{tame}_{{}^cG,F,\iota}$ 上自对偶的极大 Cohen–Macaulay 层。特别地，它在 $\Loc^{tame}_{{}^cG,F,\iota}$ 的光滑轨迹上是有限局部自由的。
\end{cor}

注意我们把 $\mathrm{CohSpr}^{unip}_{{}^cG,F,\iota}$ 看作 $\Loc^{tame}_{{}^cG,F,\iota}$ 上的凝聚层。

\begin{proof}
由命题 4.14 即得。
\end{proof}

\begin{exam}{4.17}
设 $G=\mathrm{PGL}_2$，故 $\hat G=\mathrm{SL}_2$ 且 ${}^cG=\mathrm{GL}_2$。则在 $\Lambda=\Z[1/2q(q+1)]$ 上可以证明
\[
\mathrm{CohSpr}^{unip}_{{}^cG,F,\iota}\simeq\OO_{red}\Loc^{unip}_{{}^cG,F,\iota}\oplus\OO_{\Loc^{ur}_{{}^cG,F,\iota}}.
\]
细节见 [26]。
\end{exam}

\begin{rem}{4.18}
[88] 将证明：当 $G$ 非分歧且 $\Lambda$ 是特征零的域时，猜想 4.15 成立。
\end{rem}

\ysrc{100}
我们有如下猜想。\footnote{关于这个猜想的历史，就我们所能追溯的，作如下说明。这一猜想的某种形式几年前由 Ben-Zvi、Helm 与 Nadler 作为他们此前工作的自然延续而研究过。Hellmann 在研究 $p$-进自守形式与 $p$-进伽罗瓦表示时独立提出了类似的猜想（详见他的文章 [38]）。我们是在试图把 [77] 的工作推广到 Iwahori 层次结构时得到这些想法的（讨论见 §4.7）。我们的表述强调一般系数，这反映了我们希望在框架内理解层次的提升与下降的算术含义。相当值得注意的是，不同的考虑引向了同一个对象的研究。}

\begin{conj}{4.19}
设 $G$ 在 $F$ 上拟分裂并带有固定数据，$H_I$（相应地 $H_{I(1)}$）是相应的 Iwahori（相应地 pro-$p$ Iwahori）Hecke 代数（见注 4.8）。则存在 $\Lambda$-代数的自然同构
\[
H_I\cong\big(\End_{\OO_{\Loc^{tame}_{{}^cG,F,\iota}}}\mathrm{CohSpr}^{unip}_{{}^cG,F,\iota}\big)^{\mathrm{op}},\qquad
H_{I(1)}\cong\big(\End_{\OO_{\Loc^{tame}_{{}^cG,F,\iota}}}\mathrm{CohSpr}^{tame}_{{}^cG,F,\iota}\big)^{\mathrm{op}},
\]
它们与命题 3.12 中关于不同 $\iota$ 的同构相容。特别地，有全忠实嵌入
\[
\Mod_{H_{I(1)}}\to\IndCoh(\Loc^{tame}_{{}^cG,F,\iota}),\qquad
M\mapsto\mathrm{CohSpr}^{tame}_{{}^cG,F,\iota}\otimes_{H_{I(1)}}M.
\]
此外，下列图表应当交换
\begin{center}
\begin{tikzcd}[row sep=2.4em,column sep=4em]
H_{T,I}\arrow[r,"(4.10)"]\arrow[d] & H_I\arrow[d]\\
\big(\End\mathrm{CohSpr}^{unip}_{{}^cT,F,\iota}\big)^{\mathrm{op}}\arrow[r] &
\big(\End\mathrm{CohSpr}^{unip}_{{}^cG,F,\iota}\big)^{\mathrm{op}}
\end{tikzcd}

\medskip
\begin{tikzcd}[row sep=2.4em,column sep=4em]
H_{T,I(1)}\arrow[r,"(4.10)"]\arrow[d] & H_{I(1)}\arrow[d]\\
\big(\End\mathrm{CohSpr}^{tame}_{{}^cT,F,\iota}\big)^{\mathrm{op}}\arrow[r] &
\big(\End\mathrm{CohSpr}^{tame}_{{}^cG,F,\iota}\big)^{\mathrm{op}},
\end{tikzcd}
\end{center}
其中底部映射由态射 $\Loc^{tame}_{{}^cB,F,\iota}\to\Loc^{tame}_{{}^cT,F,\iota}$ 诱导。
\end{conj}

注意在猜想中计算自同态时，$\mathrm{CohSpr}^{unip}_{{}^cG,F,\iota}$ 仍被看作 $\Loc^{tame}_{{}^cG,F,\iota}$ 上的凝聚层，与猜想 4.11 中的非分歧情形类似。

\begin{rem}{4.20}
该猜想特别蕴含应存在自然态射
\begin{equation}
\tag{4.13}
Z^{tame}_{{}^cG,F}:=H^0\Gamma(\Loc^{tame}_{G,F,\iota},\OO)\to Z(H_{I(1)}),
\end{equation}
其中 $Z(H_{I(1)})$ 是 $H_{I(1)}$ 的中心\tn{底本 (4.13) 中 $\Loc$ 的下标漏了 $c$。}，它应当使下列图表交换
\begin{equation}
\tag{4.14}
\begin{tikzcd}[row sep=2.4em,column sep=4.5em]
Z^{tame}_{{}^cG,F}\arrow[r]\arrow[d] & Z(H_{I(1)})\arrow[d,"\cong"]\\
\big(Z^{tame}_{{}^cT,F}\big)^{W_{rel}}\arrow[r,"\cong"] & \big(H_{T,I(1)}\big)^{W_{rel}}.
\end{tikzcd}
\end{equation}
\end{rem}

\ysrc{101}
这里 $T$ 表示 $G$ 的抽象 Cartan（含义例如见 [87, 1.4]），$W_{rel}$ 是 $G$ 的相对 Weyl 群。左竖直映射来自 (3.12)。（注意 $W_{rel}\cong W_{{}^cG,{}^cT}$。）右竖直同构来自 [73, 5.1]，底部同构由猜想 4.19 对驯顺分歧环面（此时 $\mathrm{CohSpr}^{tame}_{{}^cT,F,\iota}\simeq\OO_{\Loc^{tame}_{{}^cT,F,\iota}}$）的情形诱导。

我们指出，命题 4.13 的证明已经验证了非分歧环面的猜想。此外，在与 Hemo 的即将完成的工作（[88]）中，我们将证明以下结果。

\begin{thm}{4.21}
设 $\Lambda=\overline\Q_\ell$。假设 $G$ 非分歧并带有固定数据 $(B,T,e)$，且 $(U,\psi)$ 与 $I\subset K$ 如命题 4.10 中那样联系于 $(B,T,e)$。则有自然同构
\begin{equation}
\tag{4.15}
H_I\cong\End_{\OO_{\Loc^{tame}_{{}^cG,F}}}\mathrm{CohSpr}^{unip}_{{}^cG,F},
\end{equation}
它诱导全忠实嵌入
\[
\Mod_{H_I}\to\IndCoh(\Loc^{tame}_{{}^cG,F}),\qquad
M\mapsto\mathrm{CohSpr}^{unip}_{{}^cG,F}\otimes_{H_I}M.
\]
该函子把
\begin{itemize}[leftmargin=1.6em,itemsep=0.2em]
\item $H_I$ 的反球模 $M_{\mathrm{asp}}$（见命题 4.10）送到 $\OO_{\Loc^{unip}_{{}^cG,F}}$；
\item ${}_IH_K$ 送到 $\OO_{\Loc^{ur}_{{}^cG,F}}$。特别地，当 $\Lambda=\overline\Q_\ell$ 时猜想 4.11 成立。
\end{itemize}
\end{thm}

该定理实际上由下文陈述的定理 4.44 推出。我们指出，Hellmann 在此方向上已得到部分结果（见 [38]）。此外，Ben-Zvi–Chen–Helm–Nadler 在 $G$ 分裂时也证明了同构 (4.15)（[5]）。

本小节最后讨论当 $G$ 非分歧时 $\mathrm{CohSpr}^{tame}_{{}^cG,F,\iota}$ 与 $\mathrm{CohSpr}^{unip}_{{}^cG,F,\iota}$ 的关系。首先在此情形，如刚才所述，由命题 4.13（的证明），群代数 $\Lambda T(\kappa_F)\subset H_{T,I(1)}$ 作用在 $\mathrm{CohSpr}^{tame}_{{}^cG,F,\iota}$ 上。

\begin{lem}{4.22}
存在自然同构
\[
\mathrm{CohSpr}^{tame}_{{}^cG,F,\iota}\otimes_{\Lambda T(\kappa_F)}\Lambda\cong\mathrm{CohSpr}^{unip}_{{}^cG,F,\iota},
\]
其中 $\Lambda T(\kappa_F)\to\Lambda$ 是增广映射。
\end{lem}

\begin{proof}
由命题 4.13（的证明），下图中右方块是笛卡尔的
\begin{center}
\begin{tikzcd}[row sep=2.2em,column sep=2.6em]
\Loc^{unip}_{{}^cB,F,\iota}\arrow[r]\arrow[d] &
\Loc^{ur}_{{}^cT,F,\iota}\arrow[r]\arrow[d] & \{1\}\arrow[d]\\
\Loc^{tame}_{{}^cB,F,\iota}\arrow[r] &
\Loc^{tame}_{{}^cT,F,\iota}\arrow[r] &
({}^{cl}\mathcal{R}_{\kappa_{\breve F}^\times,\hat T})^\sigma.
\end{tikzcd}
\end{center}
左方块由定义也是笛卡尔的（见 (3.22)）。故
\[
\OO_{\Loc^{unip}_{{}^cB,F,\iota}}
=\OO_{\Loc^{tame}_{{}^cB,F,\iota}}\otimes_{\Lambda[({}^{cl}\mathcal{R}_{\kappa_{\breve F}^\times,\hat T})^\sigma]}\Lambda
=\OO_{\Loc^{tame}_{{}^cB,F,\iota}}\otimes_{\Lambda T(\kappa_F)}\Lambda.
\]
由于沿 $\pi^{tame}$ 的前推与余极限交换，引理得证。
\end{proof}

\ysrc{102}
\subsection{猜想中的凝聚层}
带着前两小节的猜想，自然更进一步猜测：对每个开紧子群 $K\subset G(F)$，存在 $\Loc_{{}^cG,F,\iota}$ 上的凝聚层 $\mathcal{A}_{G,K}$，其在 $\Coh(\Loc_{{}^cG,F,\iota})$ 中的（反）自同态代数 $\End\mathcal{A}_{G,K}$ 就是 $H_K$。本小节的目标是精确表述该猜想。\footnote{当 $G$ 分裂时，密切相关的猜想也出现在 [38] 中。}我们一次性固定导子为 $\mathcal{O}_F$ 的加性特征 $\psi_0:F\to\Lambda^\times$（关于对它的选取的依赖性的讨论见注 4.5）。所有叠都基变换到 $\Lambda$。

回忆注 3.8 中关于 $\Loc_{{}^cG,F,\iota}$ 上凝聚层范畴的约定。回忆该范畴的分解 (3.8)。它经由 (3.7) 受 $\mathbf{Tor}_{Z_G,\mathrm{iso}F}$ 作用，因此每个直和项受 $\mathbf{Tor}^0_{Z_G}\subset\mathbf{Tor}_{Z_G,\mathrm{iso}F}$ 作用。另一方面，$\mathbf{Tor}^0_{Z_G}$ 也经由 (4.2) 作用在 $\Rep(G(F),\Lambda)$ 上。回忆 §4.2 中的 $\mathbf{Tor}^0_{Z_G}$-挠子 $\mathbf{W}_G$（若 $G$ 拟分裂）与一般 $G$ 的 $\mathbf{Tor}_{Z_G,\mathrm{iso}F}$-挠子 $\mathbf{TS}_G$。此外回忆，若 $\breve F/F$ 驯顺，则有谱 Deligne–Lusztig 叠 (3.18) 与 (3.21)。我们还将使用以下记号。

\begin{notn}{4.23}
注意每个权 $\lambda\in\mathbb{X}^\bullet(\hat T^{\bar\tau})$ 给出 $\hat T^{\bar\tau}/\hat T$ 上的线丛，从而经拉回
\[
\widetilde{\mathrm{Loc}}^{tame}_{{}^cG,F,\iota}\xrightarrow{\ \mathrm{pr}\ }\hat B^{\bar\tau}/\hat B\to\hat T^{\bar\tau}/\hat T
\]
给出 $\widetilde{\mathrm{Loc}}^{tame}_{{}^cG,F,\iota}$ 上的线丛。我们把这线丛记作 $\OO(\lambda)$。若 $\mathcal{F}$ 是 $\widetilde{\mathrm{Loc}}^{tame}_{{}^cG,F,\iota}$ 上的（复形）凝聚层，为简单起见把 $\mathcal{F}\otimes\OO(\lambda)$ 写作 $\mathcal{F}(\lambda)$。
\end{notn}

\begin{conj}{4.24}
固定 $t\in\mathbf{TS}_G$，令 $\beta\in\mathbb{X}^\bullet(Z_{\hat G}^{\Gamma_F})$ 是由 $t$ 确定的元素。
\begin{enumerate}[label=(\arabic*),leftmargin=1.8em,itemsep=0.3em]
\item 存在 $\mathbf{Tor}^0_{Z_G}$-等变的全忠实嵌入
\[
\mathcal{A}_G:\Rep_{\mathrm{f.g.}}(G(F),\Lambda)\to\Coh^{-\beta}(\Loc_{{}^cG,F,\iota}),
\]
它与命题 3.12 中关于不同 $\iota$ 的同构相容。应当有函子的自然同构
\[
\mathcal{A}_G\circ\mathbb{D}^{\mathrm{coh}}\cong{}'\mathbb{D}^{\mathrm{Se}}\circ\mathcal{A}_G:
\Rep_{\mathrm{f.g.}}(G(F),\Lambda)\to\Coh^{-\beta}(\Loc_{{}^cG,F,\iota}),
\]
其中 $\mathbb{D}^{\mathrm{coh}}$ 来自 (4.3)，${}'\mathbb{D}^{\mathrm{Se}}$ 来自 (3.9)。
\item 诱导的保持余极限的函子
\[
\Ind\Rep(G(F),\Lambda)\to\Ind\Coh(\Loc_{{}^cG,F,\iota})
\]
仍记作 $\mathcal{A}_G$。若 $\beta=0$（于是特别地 $G$ 拟分裂并带有固定数据），则
\[
\mathcal{A}_G(\mathrm{Whit}_{U,\psi})\cong\OO_{\Loc_{{}^cG,F,\iota}},
\]
其中 $\mathrm{Whit}_{U,\psi}$ 是由该固定数据确定的 Whittaker 模（见 (4.5)）。

对 $G(F)$ 的每个开紧子群 $K$，令 $\mathcal{A}_{G,K}:=\mathcal{A}_G(\delta_K)$。则 $\mathcal{A}_{G,K}$ 应当属于 $\Coh(\Loc_{{}^cG,F,\iota})^{\heartsuit}$。令
\[
\mathcal{A}_{G,\{1\}}:=\mathcal{A}_G(\delta_{\{1\}})\simeq\mathcal{A}_G\big(\varinjlim_K\delta_K\big)=\varinjlim_K\mathcal{A}_{G,K}.
\]
则它是 $\Loc_{{}^cG,F,\iota}$ 上的通常的 ind-凝聚层，带有 $G(F)$ 的作用（因为 $\delta_{\{1\}}$ 经由左、右正则表示成为 $G(F)\times G(F)$-表示）。

\ysrc{103}
于是 $\mathcal{A}_{G,\{1\}}$ 到 $\Loc_{{}^cG,F,\iota}$ 的每个连通分支 $D$ 的限制应当是 $D$ 上有限生成的 $\OO_D[G(F)]$-模。
\item 假设 $G$ 在驯顺分歧扩张 $\breve F/F$ 上分裂。设 $\gamma$ 是 $-\beta$ 到 $\mathbb{X}^\bullet(Z_{\hat G}^{I_F})$ 的提升，并按 (4.9) 写成 $\gamma=w_\gamma\lambda_\gamma$。设 $I_\gamma$（相应地 $I_\gamma(1)$）是相应的 Iwahori（相应地 pro-$p$ Iwahori）子群。则
\[
\mathcal{A}_G(\delta_{I_\gamma(1)})\simeq\tilde\pi_*\OO_{\Loc^{tame,w_\gamma}_{{}^cG,F,\iota}}(\lambda_\gamma),\qquad
\mathcal{A}_G(\delta_{I_\gamma})\simeq\tilde\pi_*\OO_{\Loc^{unip,w_\gamma}_{{}^cG,F,\iota}}(\lambda_\gamma).
\]
若 $G=G^*$ 非分歧且 $K$ 是由 $t$ 确定的超特殊子群，则
\[
\mathcal{A}_{G,K}\simeq\OO_{\Loc^{ur}_{G,F}}.
\]
\item 设 $P\subset G$ 是有理抛物子群，$M$ 是其 Levi 商。函子 $\mathcal{A}_M$ 与 $\mathcal{A}_G$ 还应当在表示侧与抛物诱导、在谱侧与命题 3.22 的谱抛物诱导相容。
\end{enumerate}
\end{conj}

我们将在下文讨论函子 $\mathcal{A}_G$ 如何依赖 $t\in\mathbf{TS}_G$ 的选取。但先在某些情形下把猜想的 (3) 说清楚。

\begin{exam}{4.25}
设 $G=G^*$ 驯顺分歧且 $\beta=0$。取 $\gamma=0\in\mathbb{X}^\bullet(Z_{\hat G}^{I_F})$，故 $\lambda_\gamma=0$、$w_\gamma=1$。此时猜想的 (3) 说
\[
\mathcal{A}_{G,I(1)}\simeq\mathrm{CohSpr}^{tame}_{{}^cG,F,\iota},\qquad
\mathcal{A}_{G,I}\simeq\mathrm{CohSpr}^{unip}_{{}^cG,F,\iota},
\]
这与猜想 4.19 一致。此外，猜想 4.19 中预期的交换图表也与 (4) 一致。
\end{exam}

\begin{exam}{4.26}
设 $G=D^\times/F^\times$，其中 $D$ 是 $F$ 上不变为 $1/n$ 的 $n$ 次中心除代数。则 $G$ 是 $\mathrm{PGL}_n$ 的内型，故 $\hat G=\mathrm{SL}_n$。注意
\[
\gamma=-\beta=-\alpha=1/n\in\mathbb{X}^\bullet(Z_{\hat G})\cong\Z/n.
\]
设 $w=(12\cdots n)\in W=S_n$ 是循环置换。设 $\omega_i:\hat T\to\mathbb{G}_m$ 是 $\hat T$ 的第 $i$ 个基本权。则
\[
\mathcal{A}_{G,I(1)}\simeq\tilde\pi_*\OO_{\Loc^{tame,w}_{{}^cG,F,\iota}}(\omega_1),\qquad
\mathcal{A}_{G,I}\simeq\tilde\pi_*\OO_{\Loc^{unip,w}_{{}^cG,F,\iota}}(\omega_1).
\]
可以证明，当 $D$ 是 $F$ 上四元数代数、$\Lambda=\Z_\ell$ 且 $\ell>2$、$\ell\mid q-1$ 时，$\mathcal{A}_{G,I}$ 在 $\Loc^{tame}_{{}^cG,F,\iota}$ 中由平凡表示给出的点处的完备化，与 Manning 研究过的局部形变环上的模重合 [58]。更多讨论见 [26]。
\end{exam}

\begin{rem}{4.27}
猜想的 (1) 与 (2) 将蕴含 $\mathcal{A}_{G,K}$ 是极大 Cohen–Macaulay 层。若 $\beta=0$（故 $G$ 拟分裂），我们进一步猜测它关于通常的（又称非修正的）Grothendieck–Serre 对偶是自对偶的。凝聚 Springer 层的情形见推论 4.16。
\end{rem}

\begin{rem}{4.28}
令 $\Lambda=W(\overline{\mathbb{F}}_\ell)$。当 $G=\mathrm{GL}_n$ 时，层 $\mathcal{A}_{\mathrm{GL}_n,\{1\}}$ 应当同构于 Emerton–Helm 层 $\mathcal{A}_{\mathrm{EH}}$，后者在族中插值 $\mathrm{GL}_n$ 的局部 Langlands 对应（构造见 [25,38–41]，关于这一点的讨论特别见 [38]）。另一方面，受几何 Langlands 中 Braverman–Finkelberg 的一个猜想（[11]）启发，我们有
\end{rem}
